\subsection{Transport resolvents and preservation of memory}
\label{subsec:ii-memoria-resolvente}

The extended history determines the windows
\(I_m=\{9m+1,\ldots,9m+9\}\), for \(m\geq0\).
With the numbering in this section, \(I_m=E_m\) during the first cycle.
Let \(\sigma_m\in\{-1,1\}\) be the orientation preserved by that history;
for \(0\leq m<12\), it coincides with \(\Sigma_{12}(m)\). Define
\[
 a_m=\sigma_m\sum_{t\in I_m}\mathbf1_{\{2,4,6,8\}}(d_t),
 \qquad |a_m|\leq9.
\]
The register and the moving frame are extended by
\[
 z_{m+1}=z_m+a_m\mathbf e_{m\bmod12},\qquad y_m=S^{-m}z_m.
\]
The initial state \(z_0\) preserves the previous history; it is zero for
an initially empty prefix. In this subsection, write
\(R=S^{-1}\) for frame transport. Update
\eqref{eq:eta} is maintained in every window. Subsequent orientation and
events are obtained from the extended history: return of the frame
does not presuppose periodicity of the event sequence.

\begin{proposition}[Truncation identity with terminal state]
For $N\geq1$, let
\[
Y_N(u)=\sum_{m=0}^N y_mu^m,\qquad
H_{\eta,N}(u)=\sum_{m=0}^{N-1}\eta_mu^m.
\]
The exact polynomial identity holds:
\begin{equation}
(I-uR)Y_N(u)=y_0+uH_{\eta,N}(u)-u^{N+1}Ry_N.
\label{mem:identidad-finita}
\end{equation}
\end{proposition}
\begin{proof}
In degrees $1\leq m\leq N$, the coefficient on the left-hand side is
$y_m-Ry_{m-1}=\eta_{m-1}$, by \eqref{eq:eta}.
Its coefficients in degrees zero and $N+1$ are $y_0$ and $-Ry_N$,
respectively. Comparison of coefficients proves the equality.
\end{proof}

The terminal term preserves the transported memory when a history is cut.
The identity allows a finite evaluation to be compared with its extension
while keeping explicit the state left at the boundary.


\subsubsection{Correspondence with the common memory development}

Changing the window numbering does not change the register. In the
first cycle, the conventions of the two developments are related by
\[
 E^{(I)}_{m+1}=I_m=E^{(II)}_m,\qquad 0\leq m<12.
\]
The signs and events in each window coincide after this
index change. The transport convention also coincides:
\(S\mathbf e_j=\mathbf e_{j+1\bmod12}\) and \(R=S^{-1}\).
Therefore, update \eqref{eq:eta} is the same identity as
equation~\textup{(\HMTIIRefOriginal{mem:actualizacion})} in the common
core, with the same preserved previous state.

Proposition~\HMTIIRefOriginal{mem:prop-serie} proves the generating-series
identity, its convergence for \(|u|<1\), and the uniform tail
bound~\textup{(\HMTIIRefOriginal{mem:cota})}. The finite identity
\eqref{mem:identidad-finita} adds the boundary term
\(-u^{N+1}Ry_N\) here: it allows the history to be cut without replacing its terminal
state by zero. In both cases, the parameter \(u\) counts windows
of nine transitions.

Elimination of memory components preserves the operator, source,
and readout according to Proposition~\HMTIIRefOriginal{mem:prop-schur} and
equations~\textup{(\HMTIIRefOriginal{mem:schur-fuente})} and
\textup{(\HMTIIRefOriginal{mem:schur-lector})}. The example of the cycle
of twelve windows~\textup{(\HMTIIRefOriginal{mem:ejemplo-ciclo})}
distinguishes compression of the resolvent from the resolvent of the compression.
Composition of three segments and transitivity of reduction
retain their proofs alongside identities
\textup{(\HMTIIRefOriginal{mem:cociclo-tres})} and
\textup{(\HMTIIRefOriginal{mem:schur-transitivo})}.
These references belong to the complete development included in this same
article. Balance \eqref{mem:balance} and the terminal state in
\eqref{mem:identidad-finita} preserve, in their present application,
the increments and the boundary of the register; they do not by themselves
determine additional analytic coefficients.
